\section{Research Directions \& Open Questions}

\subsection{Scalable replay and prioritization}

在高维任务中，传统 replay buffer 不足以覆盖需要探索的新状态。未来研究需要同时解决:

\begin{itemize}
    \item 如何让 prioritized sampling 兼容 multi-step TD 估计。
    \item 如何在保证 sampling 多样性的前提下整合 demonstration / human corrections。
    \item 如何让 buffer adapt to policy shift（如 `exponentially weighted recency`）。
\end{itemize}

\begin{importantbox}{Open problem: replay hygiene at scale}
在分布式训练中各 worker 写入 buffer 时需确保非重采样、请勿漏采，未来可用 vector clock 来追踪 experience provenance。
\end{importantbox}

\subsection{Robust evaluation and failure diagnostics}

Q-Learning 的 failure 通常隐蔽：reward hack, return drop, policy collapse。除了 simulation metrics，研究还需解决:

\begin{itemize}
    \item 提取 `evidence matrix` 并自动关联到 actions。
    \item 构造 `mirrored failure suite` 以 replay earlier catastrophic transitions。
    \item 在 production 中保留 `checkpoints` 供 human-in-the-loop 查证。
\end{itemize}

\begin{warningbox}{诊断中最常见的盲点}
工程团队往往只关注 reward curve，一旦 reward plateau 误判为 success 就错过了 action drift。建议 combined view (reward + TD error + success rate)。
\end{warningbox}

\subsection{Agent orchestration and governance}

Q-Learning 进入 agent stack 后，治理就需要跨 levels：

\begin{itemize}
    \item 让 metric exports 透明（TD error, max Q).
    \item 让 training job 控制 plane 记录 (checkpoint id, data seed, guardrail status).
    \item 设定 `deployment kill switch` 以 `reward drop > δ` 触发 rollback.
\end{itemize}

\begin{knowledgebox}{Governance checklist}
1. 每个 release 带上 evidence matrix + action timeline；2. 训练 pipeline 必须把 `reward variance` 率先视为 safety signal；3. Deployment 过程须支持 `fast rollback`。
\end{knowledgebox}

\subsection{本章小结}

研究方向集中在 scalable replay、robust evaluation 和 agent-level governance；这些工作将决定 Q-Learning 在更复杂任务中的可信度。
最后一页提示了未来 research direction：multi-step q-learning、sim-to-real bridging 和 safety evaluation。

\begin{figure}[H]
\centering
\includegraphics[width=0.92\textwidth,page=7]{06_cs224r_qlearning_2025.pdf}
\caption{Slide：Research directions for safe Q-Learning}
\end{figure}
\newpage
\subsection{Sim-to-real & 数据迁移}

Q-Learning 部署在真实系统前常需 sim-to-real transfer：

\begin{itemize}
    \item 在 sim 中加入 observation noise 与 actuator delay 再训练 Q。
    \item 通过 behavior cloning 将 sim checkpoint 融入 target replay buffer 中。
    \item 使用 policy distillation 让 smaller net approximate sim net。
\end{itemize}

\begin{importantbox}{迁移监管要点}
1. 确认 sim 与 real timeline alignment；2. 模拟数据必须包含真实 field noise；3. rollout 全量记录供 human-in-the-loop 审查。
\end{importantbox}

